Este apêndice apresenta apenas um trecho sanitizado do arquivo
\path{transmitter/thesis-final/main.cpp}\footnote{Disponível em: \url{https://github.com/cirillom/WCAN/blob/thesis-final/transmitter/thesis-final/main.cpp}.}. O objetivo é mostrar a estrutura de inicialização, o recebimento da configuração de \emph{boot} por UART e a seleção do papel da placa na campanha experimental, evitando que o apêndice reproduza o arquivo completo.

\begin{codigo}[language=C++,basicstyle=\ttfamily\scriptsize,breaklines=false,columns=fullflexible,caption={Trecho principal do \emph{firmware} WCAN},label={cod:main-cpp}]
extern "C" void app_main(void)
{
    const wcan_test::TestConfig config =
        wcan_test::UartTestProtocol::wait_for_boot_config();
    const wcan_test::WcanTestSession test_session(config);

    if (config.role == wcan_test::Role::kIdle) {
        test_session.ready();
        test_session.wait_idle_start();
        std::printf("WCAN_TEST_END role=\%s\n", config.role_name());
        std::fflush(stdout);
        return;
    }

    init_nvs();
    init_wifi();

    std::unique_ptr<wcan::Transceiver> transceiver;
    if (config.role == wcan_test::Role::kSensor) {
        transceiver = std::make_unique<wcan::Transceiver>(
            std::vector<wcan::CANId_t>{},
            config.sensor_tx_ids(),
            config.linger_ms,
            true);
    } else if (config.role == wcan_test::Role::kReceiver) {
        transceiver = std::make_unique<wcan::Transceiver>(
            config.receiver_rx_ids(),
            std::vector<wcan::CANId_t>{},
            0,
            config.receiver_filter_enabled);
    }

    if (!transceiver || !transceiver->init()) {
        test_session.abort("setup");
        return;
    }

    std::unique_ptr<wcan_sensor::RampCanSensor> sensor;
    if (config.role == wcan_test::Role::kSensor) {
        sensor = std::make_unique<wcan_sensor::RampCanSensor>(*transceiver);
    }

    test_session.ready();
    test_session.run(*transceiver, sensor.get());
}
\end{codigo}

O código-fonte completo do projeto está disponível em
\url{https://github.com/cirillom/WCAN}.
